\subsubsection*{Feature Engineering and Selection}

The feature set was constructed through a multi-stage process combining domain-driven feature engineering with data-driven selection and refinement. The objective was to capture the temporal dynamics of soil moisture while maintaining a compact and stable representation for modeling.

\subsubsection*{Feature Engineering Overview}

Starting from the preprocessed time series, a large set of candidate features was generated by applying structured transformations to satellite, meteorological, and derived signals. These features were grouped into families based on the type of information they encode.

\begin{table}[H]
\centering
\begin{tabularx}{\textwidth}{lX}
\toprule
\textbf{Family} & \textbf{Description} \\
\midrule
\textbf{A (Change dynamics)} & Differences, gradients, and percent changes capturing short-term wetting and drying \\
\textbf{B (Volatility)} & Rolling statistics and exponential moving averages capturing variability and trend \\
\textbf{C (Memory)} & Lagged observations and exponentially weighted memory terms \\
\textbf{D (Seasonality)} & Cyclical encodings and periodic descriptors \\
\textbf{G (Meteorological)} & Rainfall accumulation, API, and time-since-rain features \\
\textbf{S (SMAP-derived)} & Soil moisture signals and temporal transforms \\
\bottomrule
\end{tabularx}
\caption{Feature families used in the engineering process.}
\end{table}

This process resulted in a high-dimensional feature space (approximately 400 features) with intentional redundancy across temporal scales. While beneficial for representation learning, this redundancy requires controlled pruning during model construction.

\subsubsection*{Feature Pool Construction}

The initial candidate set was constructed from the training split by removing the target variable and identifier columns. Features were then filtered using family-based inclusion rules, restricting the search space to engineered temporal features (families A, B, C, D, G, and S). This resulted in an initial pool of approximately 408 candidate features.

\subsubsection*{Model-Based Feature Pruning}

Feature selection was performed using a two-stage wrapper method based on gradient boosted trees (XGBoost) \cite{chen2016xgboost}. At each pruning round, feature importance was estimated using permutation importance on the validation set, averaged across multiple random seeds to improve stability \cite{fisher2019}.

Let $I_j$ denote the seed-averaged permutation importance of feature $j$. At round $t$, features were ranked by $I_j$, and the lowest-ranked fraction was removed:

\begin{equation}
F_t' = \operatorname{Top}_{n_{\text{keep},t}}(I_j)
\end{equation}

The candidate subset was accepted only if its validation $R^2$ remained within a predefined tolerance of the best score observed so far; otherwise, the update was rejected and the algorithm reverted to the previous best subset.

\begin{table}[H]
\centering
\begin{tabular}{ll}
\toprule
\textbf{Stage} & \textbf{Description} \\
\midrule
Coarse pruning & Aggressive reduction to remove clearly uninformative features \\
Fine pruning & Conservative refinement with tighter performance constraints \\
\bottomrule
\end{tabular}
\caption{Two-stage feature pruning strategy.}
\end{table}

This procedure reduced the feature set from approximately 408 candidates to a high-performing subset of roughly 100 features while maintaining or improving predictive performance.

\subsubsection*{Redundancy Reduction and Final Feature Set}

Although the pruned feature set achieved strong validation performance, it still contained substantial redundancy, primarily due to high autocorrelation between features derived from similar signals at different temporal scales.

A final refinement step was therefore applied to remove highly correlated and near-duplicate features while preserving predictive signal. This step reduces overparameterization and improves model stability.

The final base model was trained using a reduced set of 38 features, retaining the most informative temporal and physical signals identified during pruning.

\subsubsection*{Final Feature Composition}

\begin{table}[H]
\centering
\begin{tabular}{ll}
\toprule
\textbf{Category} & \textbf{Examples} \\
\midrule
SMAP-derived & Smoothed moisture, AM/PM differences, gradients \\
Meteorological & Rainfall accumulation, API-based features \\
Temporal statistics & Rolling means, extrema, exponential averages \\
Lag features & Delayed system response signals \\
Seasonality & Sine/cosine encodings of annual cycle \\
Static features & Elevation, slope, aspect, soil fractions \\
\bottomrule
\end{tabular}
\caption{Composition of the final feature set (38 features).}
\end{table}

This combination enables the model to capture both fast event-driven dynamics, such as rainfall-induced wetting, and slower processes, including seasonal trends and soil memory effects, while maintaining a manageable feature space.
